\section{Related Work}
\label{sec:related}

\paragraph{Simulated and real-world robot evaluation.}
Simulated evaluations offers perfect reproducibility, and can be efficiently parallelized. Thus, numerous simulated robot benchmarks have been proposed over the years~\citep{tassa2018deepmind, todorov2012mujoco, kolve2017ai2, james2020rlbench, lee2021ikea, mees2022calvin, srivastava2022behavior, liu2023libero, bao2023dexart, xiang2023rmbench, pumacay2024colosseum, nasiriany2024robocasa}. However, simulated environments are an imperfect replica of the real world, and even works that optimize for visual and physical fidelity~\citep{li24simpler} are often limited in the diversity of tasks they support, and sometimes fail to accurately reflect real-world policy performance~\citep{zhou2025autoeval}. In the real world, benchmarks try to ensure reproducibility through detailed manuals for reproducing environment setups~\citep{calli2015ycb, yang2019replab, heo2023furniturebench, luo2023fmb, collins2023ramp, khargonkar2024scenereplica}, enabling remote access to centrally hosted evaluations~\citep{pickem2017robotarium, bauer2022real, zhou2023train, zhou2025autoeval}, or organizing in-person challenges~\citep{kitano1997robocup, krotkov2018darpa, correll2016analysis, yenamandra2024towards, EarthRoverChallengeWebsite}. However, to ensure reproducibility of initial conditions and make centralized evaluation feasible, these approaches typically restrict evaluations to a narrow set of tasks and environments, or even a specific challenge date. %
In this work, we propose an alternative approach for real-world robot evaluation based on crowd-sourcing pairwise policy comparisons across a \emph{decentralized} network of evaluators. Our approach can cover a wider distribution of tasks and environments, and is inherently more resilient and scalable than existing, centralized evaluation approaches.

\paragraph{Generalist robot policies.} 
Recently, robotics has seen a trend towards training \emph{generalist} robot policies~\citep{driess2023palm, rt22023arxiv, open_x_embodiment_rt_x_2023, octo_2023, Doshi24-crossformer, kim2024openvla, black2024pi_0, wen2024tinyvlafastdataefficientvisionlanguageaction, zhen20243dvla, belkhale2024minivla, szot2024multimodal, pertsch2025fast, geminirobotics2025, wen2025dexvla, bjorck2025gr00t} on diverse datasets of robot experience~\citep{dasari2019robonet, ebert2021bridge, walke2023bridgedata, open_x_embodiment_rt_x_2023, khazatsky2024droid, bharadhwaj2023roboagent, fang2024rh20t, shafiullah2023bringingrobotshome, contributors2025agibotworld}. Notably, multiple works have demonstrated policies that can perform tasks across many environments \emph{out of the box}~\citep{gupta2018robotlearninghomesimproving, chi2024universal, etukuru2024robotutilitymodelsgeneral, pertsch2025fast}. At the same time, conventional approaches to robot evaluation struggle to comprehensively evaluate the performance of such generalist policies, since they are cumbersome to scale to broad distributions of tasks and environments, motivating our distributed evaluation approach in this work.

\paragraph{Crowd-sourced benchmarks.} 
In recent years, crowd-sourced benchmarks have gained popularity, from language modeling~\citep{chiang2024chatbot}, to generative image and video modeling~\citep{jiang2024genai}, or even for specialized legal model training~\citep{guha2023legalbench}. In robotics, \citet{dasari2022rb2} proposed a benchmark that aggregates policy ranking results across institutions to obtain more comprehensive evaluations of learning \emph{algorithms}, but it required retraining of policies for every new institution and task at hand, and thus remained limited to a small number of tested environments and tasks. In contrast, our work proposes a distributed benchmark for \emph{generalist robot policies}, which can be evaluated out-of-the-box, and thus makes it feasible to evaluate on a much broader distribution of tasks and environments. %

